\documentclass[aspectratio=169,10pt]{beamer}

% Shared Beamer setup for Fall 2026 course slides.
% Clean Cal Poly-inspired theme: no custom class files or uncommon packages.
\mode<presentation>{
  \usetheme{default}
  \usefonttheme{professionalfonts}
}

\definecolor{CalPolyGreen}{HTML}{154734}
\definecolor{CalPolyGold}{HTML}{BD8B13}
\definecolor{CalPolySage}{HTML}{7FA28F}
\definecolor{CalPolyMint}{HTML}{DDF1E7}
\definecolor{CalPolySand}{HTML}{A29F86}
\definecolor{CalPolyBeige}{HTML}{EFEEDE}
\definecolor{DeckInk}{HTML}{1F2933}
\definecolor{DeckMuted}{HTML}{5F6B7A}
\definecolor{DeckSoft}{HTML}{F7F7F5}

\setbeamersize{text margin left=0.55in,text margin right=0.55in}
\setbeamercolor{normal text}{fg=DeckInk,bg=white}
\setbeamercolor{structure}{fg=CalPolyGreen}
\setbeamercolor{title}{fg=white}
\setbeamercolor{frametitle}{fg=CalPolyGreen,bg=white}
\setbeamercolor{block title}{bg=CalPolySage,fg=white}
\setbeamercolor{block body}{bg=CalPolyMint,fg=black}
\setbeamercolor{alerted text}{fg=CalPolyGold}

\setbeamerfont{title}{size=\Huge,series=\bfseries}
\setbeamerfont{frametitle}{size=\Large,series=\bfseries}
\setbeamerfont{block title}{size=\normalsize,series=\bfseries}

\setbeamertemplate{navigation symbols}{}
\setbeamertemplate{itemize item}{\color{CalPolyGold}$\blacktriangleright$}
\setbeamertemplate{itemize subitem}{\color{CalPolyGreen}$\bullet$}
\setbeamertemplate{footline}{
  \hfill{\color{DeckMuted}\scriptsize\insertframenumber/\inserttotalframenumber}
  \hspace{0.18in}\vspace{0.12in}
}
\setbeamertemplate{frametitle}{
  \vspace{0.18in}
  \hspace{0.02in}\insertframetitle\par
  \vspace{0.08in}
  \noindent\makebox[\linewidth][l]{%
    {\color{CalPolyGold}\rule{0.55in}{2pt}}%
    {\color{CalPolyGreen}\rule{\dimexpr\linewidth-0.55in\relax}{2pt}}%
  }
}

\newcommand{\CourseNumber}{Course}
\newcommand{\CourseTitle}{Course Title}
\newcommand{\LectureNumber}{0}
\newcommand{\LectureTitle}{Lecture Title}
\newcommand{\LectureDate}{Date}
\newcommand{\CourseUnit}{Unit}

\newcommand{\coursenumber}[1]{\renewcommand{\CourseNumber}{#1}}
\newcommand{\coursetitle}[1]{\renewcommand{\CourseTitle}{#1}}
\newcommand{\lecturenumber}[1]{\renewcommand{\LectureNumber}{#1}}
\newcommand{\lecturetitle}[1]{\renewcommand{\LectureTitle}{#1}}
\newcommand{\lecturedate}[1]{\renewcommand{\LectureDate}{#1}}
\newcommand{\courseunit}[1]{\renewcommand{\CourseUnit}{#1}}

\author{Kirk Duran}
\institute{California Polytechnic State University, San Luis Obispo}
\date{\LectureDate}

\newcommand{\makedecktitle}{
  \begingroup
  \setbeamercolor{background canvas}{bg=CalPolyGreen}
  \begin{frame}[plain]
    \vfill
    \begin{center}
      {\usebeamerfont{title}\color{white}\LectureTitle\par}
      \vspace{0.18in}
      {\Large\color{white}Lecture \LectureNumber\par}
    \end{center}
    \vfill
    \noindent{\color{white}\rule{\linewidth}{1.2pt}}\par
    \vspace{0.08in}
    \noindent\colorbox{CalPolyGold}{\makebox[0.24\linewidth][c]{\color{white}\Large\CourseNumber}}
    \hspace{0.12in}{\color{white}\Large Kirk Duran}
  \end{frame}
  \endgroup
}

\newenvironment{keyidea}{\begin{block}{Key idea}}{\end{block}}
\newenvironment{activity}{\begin{block}{In class}}{\end{block}}
\newenvironment{cpdefinition}{\begin{block}{Definition}}{\end{block}}
\newenvironment{exampleblockcp}{
  \begingroup
  \setbeamercolor{block title}{bg=CalPolySand,fg=white}
  \setbeamercolor{block body}{bg=CalPolyBeige,fg=black}
  \begin{block}{Example}
}{
  \end{block}
  \endgroup
}

\newcommand{\imageplaceholder}[2][0.3in]{%
  \noindent\makebox[\linewidth][c]{%
    \fcolorbox{CalPolySage}{CalPolyBeige}{%
      \parbox[c][#1][c]{0.86\linewidth}{%
        \centering
        {\color{DeckMuted}\scriptsize Image placeholder: }%
        {\color{DeckInk}\scriptsize\ttfamily\detokenize{#2}}%
      }%
    }%
  }%
}

\newcommand{\sectionframe}[1]{
  \begingroup
  \setbeamercolor{background canvas}{bg=CalPolyGreen}
  \begin{frame}[plain]
    \vfill
    {\color{white}\LARGE\bfseries #1\par}
    \vspace{0.18in}
    {\color{CalPolyGold}\rule{0.22\paperwidth}{3pt}}
    {\color{white}\rule{0.62\paperwidth}{3pt}}
    \vfill
  \end{frame}
  \endgroup
}

\newcommand{\placeholdernote}{
  \begin{block}{Placeholder}
    Replace these bullets with the final lecture material, examples, and in-class activities.
  \end{block}
}


\pdfstringdefDisableCommands{\def\translate#1{#1}}

\graphicspath{{../assets/lecture-15/}}

% If an image file is not yet available, the slide displays a labeled
% placeholder using the intended filename. Place the matching file in
% ../assets/lecture-15/ to replace the placeholder automatically.
\newcommand{\lectureimage}[2][1.75in]{%
  \IfFileExists{../assets/lecture-15/#2}{%
    \includegraphics[width=\linewidth,height=#1,keepaspectratio]{#2}%
  }{%
    \fbox{%
      \parbox[c][#1][c]{0.94\linewidth}{%
        \centering
        \small\textbf{Image Placeholder}\\[0.08in]
        \scriptsize\texttt{\detokenize{#2}}
      }%
    }%
  }%
}

\newcommand{\lectureplaceholder}[2][1.75in]{%
  \fbox{%
    \parbox[c][#1][c]{0.94\linewidth}{%
      \centering
      \small\textbf{Image Placeholder}\\[0.08in]
      \scriptsize #2
    }%
  }%
}

\setbeamersize{text margin left=0.42in,text margin right=0.42in}

\coursenumber{CS 1001}
\coursetitle{Introduction to Computing}
\lecturenumber{15}
\lecturetitle{Debugging and Program Inspection}
\lecturedate{Fall 2026}
\courseunit{2. Function Development and Verification}

\begin{document}

\makedecktitle

\begin{frame}{Today}
  \begin{itemize}
    \item Distinguish testing from debugging: tests detect mismatches; debugging investigates their causes.
    \item Use print statements as a simple tracing technique and identify where that technique stops scaling well.
    \item Use the VS Code visual debugger to pause execution without inserting diagnostic output into the program.
    \item Set breakpoints and inspect variables, expressions, and active call frames.
    \item Control execution using Continue, Step Over, Step Into, and Step Out.
    \item Use a hypothesis-driven process and a red--green cycle to make targeted corrections.
  \end{itemize}

  \begin{keyidea}
    Debugging is not random code editing. It is the systematic comparison of expected execution with observed execution in order to locate the first meaningful divergence.
  \end{keyidea}
\end{frame}

\sectionframe{Part 1: From a Failing Test to a Debugging Question}

\begin{frame}{Testing and Debugging Answer Different Questions}
  \begin{columns}[T,onlytextwidth]
    \begin{column}{0.46\textwidth}
      \begin{block}{Testing}
        Does the observed behavior satisfy the expected behavior?
      \end{block}

      \begin{itemize}
        \item supplies inputs;
        \item states expectations;
        \item reports pass, failure, or error.
      \end{itemize}
    \end{column}

    \begin{column}{0.46\textwidth}
      \begin{block}{Debugging}
        Why did execution produce the observed behavior?
      \end{block}

      \begin{itemize}
        \item inspects control flow;
        \item inspects program state;
        \item narrows the location of the defect.
      \end{itemize}
    \end{column}
  \end{columns}

  \begin{keyidea}
    A failed test is evidence of a mismatch. It does not, by itself, identify the statement that caused the mismatch.
  \end{keyidea}
\end{frame}

\begin{frame}{A Symptom Is Not Necessarily the Defect}
  \begin{cpdefinition}
    A \textbf{symptom} is an observable incorrect result or failure. A \textbf{defect} is a fault in the program that contributes to that incorrect behavior.
  \end{cpdefinition}

  \begin{block}{Example}
    \texttt{def absolute\_difference(a: int, b: int) -> int:}\\
    \hspace*{0.20in}\texttt{if a > b:}\\
    \hspace*{0.40in}\texttt{return a - b}\\
    \hspace*{0.20in}\texttt{return a + b}
  \end{block}

  \begin{block}{Observed symptom}
    \texttt{absolute\_difference(4, 10)} returns \texttt{14}, but the expected value is \texttt{6}.
  \end{block}

  \begin{keyidea}
    The incorrect returned value is where we notice the problem; the faulty expression occurs earlier in the execution path.
  \end{keyidea}
\end{frame}

\begin{frame}{Debugging Is Evidence-Based Program Inspection}
  \begin{cpdefinition}
    \textbf{Debugging} is the process of locating and correcting the cause of observed incorrect program behavior.
  \end{cpdefinition}

  \begin{enumerate}
    \item Reproduce a specific failure.
    \item Predict what execution should do.
    \item Observe what execution actually does.
    \item Identify the earliest meaningful difference.
    \item Correct the responsible code.
    \item Run the tests again.
  \end{enumerate}

  \begin{keyidea}
    The purpose of a debugging tool is to increase the quality of our observations; the programmer still supplies the reasoning.
  \end{keyidea}
\end{frame}

\begin{frame}{Start from a Reproducible Case}
  \begin{block}{Implementation}
    \texttt{def absolute\_difference(a: int, b: int) -> int:}\\
    \hspace*{0.20in}\texttt{if a > b:}\\
    \hspace*{0.40in}\texttt{return a - b}\\
    \hspace*{0.20in}\texttt{return a + b}
  \end{block}

  \begin{block}{Tests}
    \texttt{assert absolute\_difference(10, 4) == 6}\quad \checkmark\\
    \texttt{assert absolute\_difference(4, 10) == 6}\quad fails
  \end{block}

  \begin{activity}
    Before inspecting the program, predict which conditional path the failing input should take and what expression should produce the final value.
  \end{activity}
\end{frame}

\sectionframe{Part 2: Print Debugging}

\begin{frame}{Print Statements Can Expose Intermediate State}
  \begin{block}{Instrumented program}
    \texttt{def final\_price(price: float, rate: float) -> float:}\\
    \hspace*{0.20in}\texttt{subtotal = price}\\
    \hspace*{0.20in}\texttt{print("subtotal:", subtotal)}\\
    \hspace*{0.20in}\texttt{tax = subtotal * rate}\\
    \hspace*{0.20in}\texttt{print("tax:", tax)}\\
    \hspace*{0.20in}\texttt{total = subtotal - tax}\\
    \hspace*{0.20in}\texttt{print("total:", total)}\\
    \hspace*{0.20in}\texttt{return total}
  \end{block}

  \begin{keyidea}
    Diagnostic prints turn selected internal values into visible output, allowing us to compare actual state with our predicted state.
  \end{keyidea}
\end{frame}

\begin{frame}{Print Debugging Requires Source-Code Changes}
  \begin{columns}[T,onlytextwidth]
    \begin{column}{0.53\textwidth}
      \begin{itemize}
        \item We must decide what to print before running the program.
        \item Adding more observations requires editing the source again.
        \item Diagnostic output can become mixed with the program's intended output.
        \item Temporary prints can be forgotten after the defect is corrected.
      \end{itemize}

      \begin{block}{Important}
        Print debugging is not inherently wrong. It is a useful technique with practical limits.
      \end{block}
    \end{column}

    \begin{column}{0.41\textwidth}
      % Intended visual: source code accumulating diagnostic print statements.
      \lectureimage[2.65in]{print-debugging-clutter.png}
    \end{column}
  \end{columns}
\end{frame}

\begin{frame}[shrink=4]{From Instrumenting the Program to Inspecting the Program}
  \begin{columns}[T,onlytextwidth]
    \begin{column}{0.45\textwidth}
      \begin{block}{Print debugging}
        Modify code to emit observations during execution.
      \end{block}
    \end{column}
    \begin{column}{0.45\textwidth}
      \begin{block}{Visual debugger}
        Pause execution and inspect runtime state through the development environment.
      \end{block}
    \end{column}
  \end{columns}

  \vspace{0.12in}
  \begin{center}
    % Intended visual: print statements on left, paused debugger UI on right.
    \lectureimage[1.85in]{prints-vs-debugger.png}
  \end{center}
\end{frame}

\sectionframe{Part 3: The VS Code Visual Debugger}

\begin{frame}[shrink=4]{A Debugger Executes the Program Under Observation}
  \begin{cpdefinition}
    A \textbf{debugger} is a development tool that can suspend a running program, expose runtime state, and provide controlled operations for resuming execution.
  \end{cpdefinition}

  \begin{columns}[T,onlytextwidth]
    \begin{column}{0.53\textwidth}
      While execution is paused, we can inspect:
      \begin{itemize}
        \item the current source location;
        \item variable bindings;
        \item active function calls;
        \item selected expressions;
        \item breakpoints controlling future pauses.
      \end{itemize}
    \end{column}

    \begin{column}{0.41\textwidth}
      % Intended visual: VS Code Run and Debug view with editor, Variables, Call Stack, and toolbar.
      \lectureimage[2.25in]{vscode-debugger-overview.png}
    \end{column}
  \end{columns}
\end{frame}

\begin{frame}{A Breakpoint Requests a Controlled Pause}
  \begin{cpdefinition}
    A \textbf{breakpoint} marks a source location at which the debugger should suspend normal execution when that executable location is reached.
  \end{cpdefinition}

  \begin{block}{Example}
    \texttt{def final\_price(price: float, rate: float) -> float:}\\
    \hspace*{0.20in}\texttt{subtotal = price}\\
    \hspace*{0.20in}\texttt{tax = subtotal * rate}\quad \textit{\small breakpoint}\\
    \hspace*{0.20in}\texttt{total = subtotal + tax}\\
    \hspace*{0.20in}\texttt{return total}
  \end{block}

  \begin{keyidea}
    When execution pauses at the breakpoint, inspect the state that exists before advancing through the suspicious computation.
  \end{keyidea}
\end{frame}

\begin{frame}{The Current Line Is a Control-Flow Location}
  \begin{columns}[T,onlytextwidth]
    \begin{column}{0.54\textwidth}
      \begin{itemize}
        \item The debugger highlights the source location at which execution is currently suspended.
        \item The highlighted line identifies what is about to be executed or the current operation being stepped through, depending on the debugger state.
        \item Variables shown in the debugger represent state at that pause point.
      \end{itemize}

      \begin{alertblock}{Tracing discipline}
        Inspect state \emph{before} stepping, predict the effect of the next statement, then step and compare.
      \end{alertblock}
    \end{column}

    \begin{column}{0.40\textwidth}
      % Intended visual: highlighted current line with arrow and a breakpoint marker.
      \lectureimage[2.60in]{breakpoint-current-line.png}
    \end{column}
  \end{columns}
\end{frame}

\begin{frame}{The Variables View Shows Bindings in the Selected Frame}
  \begin{block}{Paused inside a call}
    \texttt{def final\_price(price: float, rate: float) -> float:}\\
    \hspace*{0.20in}\texttt{subtotal = price}\\
    \hspace*{0.20in}\texttt{tax = subtotal * rate}\\
    \hspace*{0.20in}\texttt{total = subtotal + tax}
  \end{block}

  \begin{center}
    \renewcommand{\arraystretch}{1.25}
    \begin{tabular}{l|l}
      \textbf{Name} & \textbf{Current value} \\\hline
      \texttt{price} & \texttt{100.0} \\
      \texttt{rate} & \texttt{0.08} \\
      \texttt{subtotal} & \texttt{100.0} \\
      \texttt{tax} & \texttt{8.0}
    \end{tabular}
  \end{center}

  \begin{keyidea}
    The Variables view is the debugger equivalent of the binding tables we have been constructing manually.
  \end{keyidea}
\end{frame}

\begin{frame}[shrink=4]{The Call Stack Shows Active Function Calls}
  \begin{columns}[T,onlytextwidth]
    \begin{column}{0.49\textwidth}
      \begin{block}{Example execution}
        \texttt{main code}\\
        \hspace*{0.16in}$\downarrow$ \texttt{final\_price(...)}\\
        \hspace*{0.32in}$\downarrow$ \texttt{compute\_tax(...)}
      \end{block}

      While \texttt{compute\_tax} is active, the earlier call to \texttt{final\_price} has not finished.
    \end{column}

    \begin{column}{0.45\textwidth}
      % Intended visual: stack frames for compute_tax, final_price, and module code.
      \lectureimage[2.20in]{call-stack-frames.png}
    \end{column}
  \end{columns}

  \begin{keyidea}
    Selecting a different frame lets us inspect the bindings associated with that active call.
  \end{keyidea}
\end{frame}

\begin{frame}{Watch Expressions Track Questions About State}
  \begin{cpdefinition}
    A \textbf{watch expression} is an expression the debugger reevaluates while execution is paused so the programmer can monitor its current value.
  \end{cpdefinition}

  \begin{block}{Examples}
    \texttt{subtotal * rate}\\
    \texttt{a > b}\\
    \texttt{result == expected}
  \end{block}

  \begin{itemize}
    \item Watches are useful when a derived value matters but is not stored in its own variable.
    \item Prefer simple observation expressions while learning; a debugger is most useful when it does not change the state being investigated.
  \end{itemize}
\end{frame}

\begin{frame}{Step Over Executes the Current Line at the Current Abstraction Level}
  \begin{block}{Paused here}
    \texttt{tax = compute\_tax(subtotal, rate)}
  \end{block}

  \begin{columns}[T,onlytextwidth]
    \begin{column}{0.46\textwidth}
      \begin{block}{Step Over}
        Execute the line, including the call to \texttt{compute\_tax}, but do not pause on each statement inside that function.
      \end{block}
    \end{column}

    \begin{column}{0.46\textwidth}
      \begin{block}{Afterward}
        Pause at the next statement in the current frame with \texttt{tax} now bound to the returned value.
      \end{block}
    \end{column}
  \end{columns}

  \begin{keyidea}
    Step Over is appropriate when the called function is not currently the object of investigation.
  \end{keyidea}
\end{frame}

\begin{frame}{Step Into Follows a Function Call}
  \begin{block}{Paused here}
    \texttt{tax = compute\_tax(subtotal, rate)}
  \end{block}

  \begin{columns}[T,onlytextwidth]
    \begin{column}{0.46\textwidth}
      \begin{block}{Step Into}
        Enter the called function and suspend execution inside its body.
      \end{block}
    \end{column}

    \begin{column}{0.46\textwidth}
      \begin{block}{New frame}
        Inspect parameter bindings and local variables belonging to the \texttt{compute\_tax} call.
      \end{block}
    \end{column}
  \end{columns}

  \begin{keyidea}
    Step Into crosses a function abstraction boundary so we can inspect the implementation of the call.
  \end{keyidea}
\end{frame}

\begin{frame}{The Four Core Execution Controls Serve Different Purposes}
  \begin{center}
    \renewcommand{\arraystretch}{1.35}
    \begin{tabular}{l|p{0.68\textwidth}}
      \textbf{Action} & \textbf{Effect} \\\hline
      Continue & Run until the next stopping event. \\
      Step Over & Execute the current line without tracing into calls on that line. \\
      Step Into & Follow a called function and inspect its execution. \\
      Step Out & Finish the current function and pause in its caller.
    \end{tabular}
  \end{center}

  \vspace{0.10in}
  % Intended visual: VS Code debug toolbar labeled with the four execution controls.
  \lectureimage[1.55in]{debug-step-actions.png}
\end{frame}

\sectionframe{Part 4: Debugging as a Controlled Trace}

\begin{frame}{The Debugger Extends the Manual Trace Model}
  \begin{columns}[T,onlytextwidth]
    \begin{column}{0.45\textwidth}
      \begin{block}{Manual trace}
        \begin{itemize}
          \item current statement;
          \item current bindings;
          \item active frame;
          \item predicted next state.
        \end{itemize}
      \end{block}
    \end{column}

    \begin{column}{0.45\textwidth}
      \begin{block}{Debugger}
        \begin{itemize}
          \item highlighted source location;
          \item Variables view;
          \item Call Stack view;
          \item controlled stepping.
        \end{itemize}
      \end{block}
    \end{column}
  \end{columns}

  \begin{keyidea}
    The debugger does not replace tracing. It lets us compare our trace model with the actual execution state one step at a time.
  \end{keyidea}
\end{frame}

\begin{frame}[shrink=5]{Tracing Example: Find the First Incorrect State}
  \begin{block}{Buggy function}
    \texttt{def final\_price(price: float, rate: float) -> float:}\\
    \hspace*{0.20in}\texttt{subtotal = price}\\
    \hspace*{0.20in}\texttt{tax = subtotal * rate}\\
    \hspace*{0.20in}\texttt{total = subtotal - tax}\\
    \hspace*{0.20in}\texttt{return total}
  \end{block}

  For \texttt{final\_price(100.0, 0.08)}, suppose the specification expects \texttt{108.0}.

  \begin{center}
    \renewcommand{\arraystretch}{1.20}
    \begin{tabular}{l|l|l}
      \textbf{Pause point} & \textbf{Expected state} & \textbf{Observed state} \\\hline
      after \texttt{subtotal} & \texttt{subtotal = 100.0} & \texttt{100.0} \\
      after \texttt{tax} & \texttt{tax = 8.0} & \texttt{8.0} \\
      after \texttt{total} & \texttt{total = 108.0} & \texttt{92.0}
    \end{tabular}
  \end{center}

  \begin{keyidea}
    The first state divergence occurs when \texttt{total} is computed, sharply narrowing the defect location.
  \end{keyidea}
\end{frame}

\begin{frame}{Place Breakpoints Near a Hypothesis, Not Everywhere}
  \begin{itemize}
    \item A breakpoint before every line produces many pauses but little prioritization.
    \item A useful breakpoint is placed near the earliest location where the failing behavior could plausibly emerge.
    \item Begin before the suspicious computation so that the inputs to that computation can also be inspected.
  \end{itemize}

  \begin{block}{Example hypothesis}
    ``The tax amount is correct, but the final price is combining subtotal and tax incorrectly.''
  \end{block}

  \begin{block}{Breakpoint choice}
    Pause on or immediately before \texttt{total = subtotal - tax}.
  \end{block}
\end{frame}

\begin{frame}{Step Over When the Callee Is Not Suspect}
  \begin{block}{Program}
    \texttt{def compute\_tax(amount: float, rate: float) -> float:}\\
    \hspace*{0.20in}\texttt{return amount * rate}\\[0.08in]
    \texttt{def final\_price(price: float, rate: float) -> float:}\\
    \hspace*{0.20in}\texttt{tax = compute\_tax(price, rate)}\\
    \hspace*{0.20in}\texttt{return price - tax}
  \end{block}

  If previous tests already establish that \texttt{compute\_tax} behaves correctly, Step Over the call and inspect the returned value in \texttt{tax}.

  \begin{keyidea}
    Debugging is selective. Do not inspect an implementation merely because the debugger can enter it.
  \end{keyidea}
\end{frame}

\begin{frame}{Step Into When the Returned Value Is Already Wrong}
  \begin{block}{Suppose}
    After stepping over \texttt{tax = compute\_tax(price, rate)}, the debugger shows \texttt{tax = -8.0} when \texttt{8.0} was expected.
  \end{block}

  \begin{enumerate}
    \item Restart the failing case.
    \item Pause before the call to \texttt{compute\_tax}.
    \item Verify the argument values.
    \item Step Into the call.
    \item Trace the callee's parameter bindings and expression evaluation.
  \end{enumerate}

  \begin{keyidea}
    Step Into is justified by evidence that the incorrect state exists at or before the function boundary.
  \end{keyidea}
\end{frame}

\begin{frame}{The Call Stack Answers ``How Did We Get Here?''}
  \begin{block}{Paused inside \texttt{compute\_tax}}
    \centering
    \texttt{compute\_tax}\quad current frame\\
    $\uparrow$\\
    \texttt{final\_price}\quad caller\\
    $\uparrow$\\
    \texttt{test\_final\_price}\quad earlier caller
  \end{block}

  \begin{itemize}
    \item Select \texttt{compute\_tax} to inspect its parameters and locals.
    \item Select \texttt{final\_price} to inspect the still-active caller frame.
    \item The call stack reconstructs the active chain of calls leading to the current location.
  \end{itemize}
\end{frame}

\begin{frame}{Conditionals Should Be Debugged by Following One Concrete Path}
  \begin{block}{Buggy function}
    \texttt{def shipping\_cost(weight: float) -> float:}\\
    \hspace*{0.20in}\texttt{if weight < 10:}\\
    \hspace*{0.40in}\texttt{return 5.0}\\
    \hspace*{0.20in}\texttt{return 8.0}
  \end{block}

  Suppose the specification says packages weighing \emph{at most} 10 pounds cost \texttt{5.0}.

  \begin{activity}
    For the failing input \texttt{10.0}, identify the expected condition value, the observed condition value, and the first point where the execution path diverges.
  \end{activity}

  \begin{keyidea}
    A conditional defect may produce the wrong branch even when every statement inside each branch is individually correct.
  \end{keyidea}
\end{frame}

\begin{frame}{The First Divergence Principle Narrows the Search}
  \begin{columns}[T,onlytextwidth]
    \begin{column}{0.54\textwidth}
      When debugging a concrete failing case:
      \begin{enumerate}
        \item predict the next relevant state;
        \item observe the actual state;
        \item advance only far enough to compare again;
        \item stop broadening the search once the first incorrect state is found.
      \end{enumerate}

      Later incorrect values may be consequences rather than independent defects.
    \end{column}

    \begin{column}{0.40\textwidth}
      % Intended visual: two execution traces that match, then diverge at one statement.
      \lectureimage[2.65in]{first-divergence.png}
    \end{column}
  \end{columns}
\end{frame}

\sectionframe{Part 5: Red--Green Development}

\begin{frame}{A Failing Test Gives Us the Red State}
  \begin{block}{RED}
    At least one test representing required behavior is currently failing.
  \end{block}

  \begin{itemize}
    \item Red is not merely ``the code is bad.''
    \item It identifies a specific observable behavior that the current implementation does not satisfy.
    \item The failure gives us a reproducible case from which debugging can begin.
  \end{itemize}

  \begin{keyidea}
    Do not erase the evidence by weakening the test simply to make it pass. First determine whether the test's expectation correctly represents the specification.
  \end{keyidea}
\end{frame}

\begin{frame}{Debugging Should Produce a Targeted Change}
  \begin{block}{Before}
    \texttt{total = subtotal - tax}
  \end{block}

  \begin{block}{Evidence}
    \texttt{subtotal = 100.0}, \texttt{tax = 8.0}, expected \texttt{total = 108.0}, observed \texttt{92.0}
  \end{block}

  \begin{block}{Targeted correction}
    \texttt{total = subtotal + tax}
  \end{block}

  \begin{keyidea}
    The edit should follow from the diagnosis. Random changes make it difficult to know which change affected the behavior and why.
  \end{keyidea}
\end{frame}

\begin{frame}{Green Means the Relevant Tests Pass}
  \begin{block}{GREEN}
    The corrected implementation now satisfies the current test suite.
  \end{block}

  \begin{itemize}
    \item First rerun the test that exposed the defect.
    \item Then run the complete relevant test suite.
    \item A local correction can accidentally change behavior elsewhere.
  \end{itemize}

  \begin{keyidea}
    Re-running existing tests after a change checks for regressions: previously expected behavior that no longer works after the modification.
  \end{keyidea}
\end{frame}

\begin{frame}[shrink=4]{Red and Green Form a Development Feedback Loop}
  \begin{center}
    % Intended visual: RED failing test -> inspect/debug -> change code -> GREEN passing suite -> next behavior.
    \lectureimage[1.85in]{red-green-cycle.png}
  \end{center}

  \begin{block}{Core loop}
    \centering
    failing expectation $\rightarrow$ inspect execution $\rightarrow$ targeted change $\rightarrow$ passing tests
  \end{block}

  \begin{keyidea}
    The cycle is useful because each code change has immediate executable feedback about whether the expected behavior was restored.
  \end{keyidea}
\end{frame}

\sectionframe{Part 6: An Integrated Debugging Workflow}

\begin{frame}{A Practical Debugging Procedure}
  \begin{enumerate}
    \item Run the tests and select one reproducible failure.
    \item Read the expected and actual values carefully.
    \item Predict the execution path for the failing input.
    \item Form a specific hypothesis about where the first divergence may occur.
    \item Set a breakpoint before that location.
    \item Inspect state, then Step Over or Step Into deliberately.
    \item Correct the defect supported by the evidence.
    \item Rerun the failing test, then the full relevant suite.
  \end{enumerate}
\end{frame}

\begin{frame}[shrink=6]{Integrated Example: Debugging Across a Function Call}
  \begin{block}{Implementation}
    \texttt{def apply\_discount(price: float, rate: float) -> float:}\\
    \hspace*{0.20in}\texttt{return price + price * rate}\\[0.05in]
    \texttt{def sale\_price(price: float, rate: float) -> float:}\\
    \hspace*{0.20in}\texttt{discounted = apply\_discount(price, rate)}\\
    \hspace*{0.20in}\texttt{return discounted}
  \end{block}

  \begin{block}{Failing case}
    \texttt{sale\_price(100.0, 0.20)} should return \texttt{80.0}, but returns \texttt{120.0}.
  \end{block}

  \begin{activity}
    Where would you place the first breakpoint? Would you initially Step Over or Step Into \texttt{apply\_discount}? What observation would justify changing that decision?
  \end{activity}
\end{frame}

\begin{frame}{Common Debugging Mistakes}
  \begin{itemize}
    \item Editing several unrelated statements before rerunning the test.
    \item Stepping through every line without a hypothesis about what information is needed.
    \item Looking only at the final incorrect value rather than finding where state first became incorrect.
    \item Stepping into every function call regardless of whether evidence implicates that function.
    \item Fixing the one failing case but not rerunning the rest of the tests.
    \item Leaving diagnostic code in place after it no longer serves a purpose.
  \end{itemize}

  \begin{keyidea}
    Effective debugging reduces uncertainty systematically; ineffective debugging often creates more uncontrolled changes than evidence.
  \end{keyidea}
\end{frame}

\begin{frame}[shrink=6]{Final Trace: Locate the Defect Before Editing}
  \begin{block}{Program}
    \texttt{def is\_eligible(age: int) -> bool:}\\
    \hspace*{0.20in}\texttt{return age > 18}\\[0.06in]
    \texttt{def admission(age: int, has\_ticket: bool) -> str:}\\
    \hspace*{0.20in}\texttt{if is\_eligible(age) and has\_ticket:}\\
    \hspace*{0.40in}\texttt{return "admit"}\\
    \hspace*{0.20in}\texttt{return "deny"}
  \end{block}

  Specification: age \texttt{18} is eligible when the person also has a ticket.

  \begin{activity}
    For \texttt{admission(18, True)}, identify: (1) the expected result, (2) the first useful breakpoint, (3) where Step Into becomes useful, and (4) the first observed state or expression that would differ from the specification.
  \end{activity}
\end{frame}

\begin{frame}{Takeaways}
  \begin{itemize}
    \item Tests identify reproducible mismatches between expected and actual behavior.
    \item Print statements can reveal state, but visual debugging scales better when detailed runtime inspection is needed.
    \item Breakpoints pause execution so variables, expressions, and active call frames can be inspected directly.
    \item Continue, Step Over, Step Into, and Step Out intentionally control the level and amount of observation.
    \item Debugging should seek the first meaningful divergence between predicted and observed execution.
    \item Red--green development connects failing tests, evidence-based corrections, and repeated verification.
  \end{itemize}

  \begin{keyidea}
    Tests tell us \emph{that} behavior is wrong. Debugging helps us explain \emph{why} it is wrong.
  \end{keyidea}
\end{frame}

\end{document}
